% GENERATED FILE. Edit canonical lessons, then run npm run generate:books.
% canonical-source-hash: fnv1a64:968cf5a741bb1e58
% canonical-lessons: GE-C443-feinkostgeschaeft, GE-C443-konditorei, GE-C443-gemuesehaendler, GE-C443-drogerie, GE-C443-reformhaus

\chapter{Kleine Läden}
\label{ch:ge-kleine-laden}
\hlchaptermodality{443}

\begin{chapteropening}
\textbf{By the end of this chapter:} \emph{I can find small food and health shops: the delicatessen, the cake shop, the greengrocer, the chemist and the health food shop.}

Complete the last lesson of chapter 443: I can find small food and health shops: the delicatessen, the cake shop, the greengrocer, the chemist and the health food shop.
\end{chapteropening}

\section[das Feinkostgeschäft]{das Feinkostgeschäft — the delicatessen}

\label{lesson:GE-C443-feinkostgeschaeft}



\begin{quote}
Before the new one: say the German for the snowplough, then the German for to slip.
\end{quote}

\subsection*{You'll want to know: das Feinkostgeschäft}
\textbf{das Feinkostgeschäft} — \textquotedblleft{}the delicatessen\textquotedblright{}.

\begin{quote}
Im Feinkostgeschäft kauft er Oliven und Schinken.
\end{quote}

\begin{quote}
\emph{At the delicatessen he buys olives and ham.}
\end{quote}

Plural: \emph{die Feinkostgeschäfte}. \emph{Die Feinkost} — fine foods.

\begin{cousinweb}
\emph{Feinkostgeschäft} is built from \emph{fein} 'fine', which came into Middle High German as \emph{fīn} from Old French \emph{fin}, \emph{Kost} 'food, fare', and \emph{Geschäft} 'business, shop', inherited from Old High German \emph{giscaft} (EtymDB).
\end{cousinweb}

\begin{grammarlens}[title={What you've built}]
One more word for telling, explaining and describing at B1.
\end{grammarlens}

\subsection*{Your turn}
\begin{itemize}
\raggedright
  \item \emph{Say it:} \emph{das Feinkostgeschäft}
  \item \emph{Say it:} \emph{das Feinkostgeschäft}, once more
  \item \emph{Say it:} a sentence of your own with \emph{das Feinkostgeschäft}
  \item \emph{Recall:} say the German for the snowplough, then the German for to slip, then say \emph{das Feinkostgeschäft} again
\end{itemize}

\begin{tcolorbox}[breakable,colback=teal!4,colframe=teal!35!black,title={Before you move on}]
What does das Feinkostgeschäft mean? (\textquotedblleft{}The delicatessen\textquotedblright{}.) Say it once more.
\end{tcolorbox}
\section[die Konditorei]{die Konditorei — the cake shop, the patisserie}

\label{lesson:GE-C443-konditorei}



\begin{quote}
Before the new one: say the German for to slip, then the German for the delicatessen.
\end{quote}

\subsection*{You'll want to know: die Konditorei}
\textbf{die Konditorei} — \textquotedblleft{}the cake shop, the patisserie\textquotedblright{}.

\begin{quote}
In der Konditorei gibt es den besten Kuchen der Stadt.
\end{quote}

\begin{quote}
\emph{The cake shop has the best cake in town.}
\end{quote}

Plural: \emph{die Konditoreien}. \emph{Der Konditor} — the pastry cook.

\begin{cousinweb}
\emph{Konditorei} is built on \emph{Konditor} 'confectioner' with the ending \emph{-ei} (EtymDB), a learned borrowing from Latin \emph{condītor} 'a seasoner', from \emph{condīre} 'to season, to preserve' (Lewis \& Short).
\end{cousinweb}

\begin{grammarlens}[title={What you've built}]
One more word for telling, explaining and describing at B1.
\end{grammarlens}

\subsection*{Your turn}
\begin{itemize}
\raggedright
  \item \emph{Say it:} \emph{die Konditorei}
  \item \emph{Say it:} \emph{die Konditorei}, once more
  \item \emph{Say it:} a sentence of your own with \emph{die Konditorei}
  \item \emph{Recall:} say the German for to slip, then the German for the delicatessen, then say \emph{die Konditorei} again
\end{itemize}

\begin{tcolorbox}[breakable,colback=teal!4,colframe=teal!35!black,title={Before you move on}]
What does die Konditorei mean? (\textquotedblleft{}The cake shop, the patisserie\textquotedblright{}.) Say it once more.
\end{tcolorbox}
\section[der Gemüsehändler]{der Gemüsehändler — the greengrocer}

\label{lesson:GE-C443-gemuesehaendler}



\begin{quote}
Before the new one: say the German for the delicatessen, then the German for the cake shop.
\end{quote}

\subsection*{You'll want to know: der Gemüsehändler}
\textbf{der Gemüsehändler} — \textquotedblleft{}the greengrocer\textquotedblright{}.

\begin{quote}
Der Gemüsehändler hat heute frische Tomaten.
\end{quote}

\begin{quote}
\emph{The greengrocer has fresh tomatoes today.}
\end{quote}

Same form in the plural. A woman: \emph{die Gemüsehändlerin}.

\begin{cousinweb}
\emph{Gemüsehändler} is built from \emph{Gemüse} 'vegetables', from Middle High German \emph{gemüese}, built on \emph{Mus} 'mash, porridge' (EtymDB), and \emph{Händler} 'dealer', from \emph{handeln} 'to trade'.
\end{cousinweb}

\begin{grammarlens}[title={What you've built}]
One more word for telling, explaining and describing at B1.
\end{grammarlens}

\subsection*{Your turn}
\begin{itemize}
\raggedright
  \item \emph{Say it:} \emph{der Gemüsehändler}
  \item \emph{Say it:} \emph{der Gemüsehändler}, once more
  \item \emph{Say it:} a sentence of your own with \emph{der Gemüsehändler}
  \item \emph{Recall:} say the German for the delicatessen, then the German for the cake shop, then say \emph{der Gemüsehändler} again
\end{itemize}

\begin{tcolorbox}[breakable,colback=teal!4,colframe=teal!35!black,title={Before you move on}]
What does der Gemüsehändler mean? (\textquotedblleft{}The greengrocer\textquotedblright{}.) Say it once more.
\end{tcolorbox}
\section[die Drogerie]{die Drogerie — the chemist (for toiletries), the drugstore}

\label{lesson:GE-C443-drogerie}



\begin{quote}
Before the new one: say the German for the cake shop, then the German for the greengrocer.
\end{quote}

\subsection*{You'll want to know: die Drogerie}
\textbf{die Drogerie} — \textquotedblleft{}the chemist (for toiletries), the drugstore\textquotedblright{}.

\begin{quote}
In der Drogerie kauft sie Shampoo und Zahnpasta.
\end{quote}

\begin{quote}
\emph{At the chemist she buys shampoo and toothpaste.}
\end{quote}

Plural: \emph{die Drogerien}. Medicine is sold at \emph{die Apotheke}.

\begin{cousinweb}
\emph{Drogerie} was borrowed from French \emph{droguerie} (EtymDB), built on \emph{drogue} 'dried goods, drug'; the deeper origin of \emph{drogue} is not settled.
\end{cousinweb}

\begin{grammarlens}[title={What you've built}]
One more word for telling, explaining and describing at B1.
\end{grammarlens}

\subsection*{Your turn}
\begin{itemize}
\raggedright
  \item \emph{Say it:} \emph{die Drogerie}
  \item \emph{Say it:} \emph{die Drogerie}, once more
  \item \emph{Say it:} a sentence of your own with \emph{die Drogerie}
  \item \emph{Recall:} say the German for the cake shop, then the German for the greengrocer, then say \emph{die Drogerie} again
\end{itemize}

\begin{tcolorbox}[breakable,colback=teal!4,colframe=teal!35!black,title={Before you move on}]
What does die Drogerie mean? (\textquotedblleft{}The chemist (for toiletries), the drugstore\textquotedblright{}.) Say it once more.
\end{tcolorbox}
\section[das Reformhaus]{das Reformhaus — the health food shop}

\label{lesson:GE-C443-reformhaus}



\begin{quote}
Before the new one: say the German for the greengrocer, then the German for the chemist.
\end{quote}

\subsection*{You'll want to know: das Reformhaus}
\textbf{das Reformhaus} — \textquotedblleft{}the health food shop\textquotedblright{}.

\begin{quote}
Das Brot aus dem Reformhaus ist teuer, aber gut.
\end{quote}

\begin{quote}
\emph{The bread from the health food shop is expensive but good.}
\end{quote}

Plural: \emph{die Reformhäuser}.

\begin{cousinweb}
\emph{Reformhaus} is built from \emph{Reform} 'reform', borrowed from French \emph{réforme} (EtymDB), which goes back to Latin \emph{reformāre} 'to reshape', and \emph{Haus} 'house', inherited from Proto-Germanic \emph{hūsą} (EtymDB). The shops grew out of a movement for a simpler, healthier life.
\end{cousinweb}

\begin{grammarlens}[title={What you've built}]
One more word for telling, explaining and describing at B1.
\end{grammarlens}

\subsection*{Your turn}
\begin{itemize}
\raggedright
  \item \emph{Say it:} \emph{das Reformhaus}
  \item \emph{Say it:} \emph{das Reformhaus}, once more
  \item \emph{Say it:} a sentence of your own with \emph{das Reformhaus}
  \item \emph{Recall:} say the German for the greengrocer, then the German for the chemist, then say \emph{das Reformhaus} again
\end{itemize}

\begin{tcolorbox}[breakable,colback=teal!4,colframe=teal!35!black,title={Before you move on}]
What does das Reformhaus mean? (\textquotedblleft{}The health food shop\textquotedblright{}.) Say it once more.
\end{tcolorbox}
